\documentclass[11pt]{article}

\usepackage[margin=1in]{geometry}
\usepackage{amsmath,amssymb}
\usepackage{booktabs}
\usepackage{array}
\usepackage{longtable}
\usepackage{xcolor}
\usepackage{listings}
\usepackage{hyperref}
\usepackage{enumitem}

\hypersetup{
  colorlinks=true,
  linkcolor=black,
  urlcolor=blue!60!black,
  citecolor=black,
  pdftitle={Bounded Analysis: scrontab Master-Runner Consolidation in generate\_workflows.sh},
  pdfauthor={EIB MCP-RAG analysis}
}

\lstset{
  basicstyle=\ttfamily\small,
  breaklines=true,
  columns=fullflexible,
  frame=single,
  framesep=4pt,
  backgroundcolor=\color{gray!8},
  showstringspaces=false
}

\newcommand{\code}[1]{\texttt{#1}}

\title{Bounded Analysis of the \code{scrontab} Master-Runner\\
       Consolidation in \code{generate\_workflows.sh}}
\author{NOAA EIB --- Global Workflow PR Test Harness Review}
\date{April~30, 2026}

\begin{document}
\maketitle

\begin{abstract}
This note evaluates the change introduced on branch
\code{copilot/modify-scrontab-scripts} of the
\code{global-workflow.DavidH} fork, which collapses $N$ per-experiment
\code{scrontab} entries into a single sequential master-runner job.
We frame the change in terms of the only function the cron/scron entry
performs --- driving the Rocoto workflow tick --- and place bounded
estimates on its memory, CPU, and throughput consequences. The
conclusion is that the absolute resource saving is well under one
percent of either the scron service node or the downstream PR-sweep
compute footprint, but is materially useful in the narrow case of a
small \code{cgroup} on the scron job (the actual Gaea failure mode).
\end{abstract}

\section{Scope and System Context}
\label{sec:scope}

The CM for \code{global-workflow} maintains an out-of-band PR test
harness driven by \code{dev/workflow/generate\_workflows.sh}. The
script generates an experiment per PR YAML case under
\code{dev/ci/cases/pr/}, and registers a scheduled tick that advances
each experiment's Rocoto workflow.

The scheduled tick has exactly one purpose: invoke
\code{rocotorun} (or the experiment's \code{*.scron.sh} wrapper) so
that Rocoto reads its XML, queries/updates the per-experiment SQLite
workflow database, and submits whatever cycle tasks have become
eligible to Slurm. The tick performs no science, holds no compute
resources, and is independent of every actual model job that Rocoto
subsequently submits to the batch system.

The branch under review changes \emph{only} how the tick is scheduled
on hosts that use Slurm-managed \code{scrontab} (Gaea). All
downstream behaviour --- per-experiment XML, SQLite database,
\code{EXPDIR}/\code{COMROOT}/\code{DATAROOT}, task graph, and resource
footprint --- is unchanged.

\section{Pre- and Post-Change Behaviour}
\label{sec:behaviour}

\paragraph{Before.}
For each experiment, the per-experiment crontab block
(\code{\${pslot}.crontab}, containing its \code{\#SCRON --*} directives
and a \code{*/5 * * * *} line invoking \code{\${pslot}.scron.sh}) was
appended to \code{tests.cron}. Installation via \code{scrontab}
created $N$ independent scron jobs --- one per experiment --- each
firing every five minutes and each launching its own \code{rocotorun}
on the scron service node concurrently.

\paragraph{After.}
The script collects each experiment's \code{*.scron.sh} path into
\code{\_scron\_sh\_files}, then writes a single
\code{rocoto\_master\_run.sh} under \code{\${RUNTESTS}/EXPDIR/} which
calls each scron script sequentially. A single
\code{master\_scron} block is emitted to \code{tests.cron} reusing the
\code{\#SCRON --partition=}/\code{--account=} directives drawn from the
first experiment's crontab and adding
\code{\#SCRON --dependency=singleton}, ensuring the next 5-minute tick
cannot start until the previous one drains.

The transformation is purely a concurrency-flattening of the tick
work: total CPU-seconds and total I/O are unchanged; only the
concurrency profile differs.

\section{Bounded Cost Model}
\label{sec:bounds}

Let $N$ be the number of PR cases, $T = 300$\,s the tick period, $R$
the resident set size of one \code{rocotorun} pass, and $t$ its
wallclock duration. Plausible ranges drawn from typical Rocoto
deployments and the GW PR XML sizes are summarised in
Table~\ref{tab:bounds}.

\begin{table}[h]
\centering
\caption{Bounding parameters per \code{rocotorun} invocation.}
\label{tab:bounds}
\begin{tabular}{lll}
\toprule
Quantity & Plausible range & Source of bound \\
\midrule
RSS of one \code{rocotorun} ($R$) & $50$--$300$\,MB &
  Ruby + XML DOM + libsqlite + \code{squeue} fork \\
Wallclock of one pass ($t$)       & $2$--$30$\,s &
  Dominated by \code{squeue}/\code{sacct} round-trip + XML parse \\
CPU-time per pass                  & $\approx t$ &
  Single-threaded Ruby \\
PR cases ($N$)                    & $\approx 10$ &
  Current sweep size; up to $\sim 30$ plausible \\
Tick period ($T$)                 & $300$\,s &
  Hardcoded \code{*/5 * * * *} \\
\bottomrule
\end{tabular}
\end{table}

\subsection{Peak Concurrent Footprint on the scron Node}

\paragraph{Before --- $N$ parallel ticks.}
\begin{align}
M_\text{peak} &= N \cdot R \;\in\; [10 \cdot 50,\; 10 \cdot 300]\,\text{MB}
              \;=\; \mathbf{0.5\text{--}3\,\text{GB}} \\
C_\text{peak} &= N \text{ cores busy for } t\,\text{s}
              \;\in\; [20,\; 300]\,\text{CPU-s per } 300\text{-s window}
\end{align}

\paragraph{After --- one serial tick.}
\begin{align}
M_\text{peak} &= R \;\in\; \mathbf{0.05\text{--}0.3\,\text{GB}} \\
C_\text{peak} &= 1 \text{ core busy for } N\cdot t\,\text{s}
\end{align}

Total CPU-seconds and total scheduler RPC count are identical.

\subsection{Duty Cycle of the Master Tick}

\begin{equation}
D \;=\; \frac{N \cdot t}{T} \;\in\;
  \left[\tfrac{10\cdot 2}{300},\; \tfrac{10\cdot 30}{300}\right]
  \;=\; [0.7\%,\; 100\%]
\label{eq:duty}
\end{equation}

The lower end is comfortably idle. The upper end --- long passes
combined with many cases --- is precisely where the
\code{singleton} dependency starts \emph{stretching} the per-experiment
polling cadence and the implicit throughput cost becomes real.
Realistic GW PR sweeps sit near $D \approx 5\text{--}15\%$, so the
nominal 5-minute cadence is preserved.

\section{What the Saving Is Worth}
\label{sec:value}

Comparing the saving against the things that actually share the
service node and against the PR sweep's downstream footprint:

\begin{table}[h]
\centering
\caption{Magnitude comparison (typical service-node and PR-sweep scales).}
\label{tab:magnitude}
\begin{tabular}{ll}
\toprule
Reference & Order of magnitude \\
\midrule
Memory saved (this PR), peak                 & $\sim 0.5\text{--}3$\,GB \\
Memory saved (this PR), average              & $\sim 50\text{--}300$\,MB \\
Typical scron-node total RAM (Gaea es-login) & $256\text{--}500$\,GB \\
\textbf{Saved fraction of node memory}       & $\boldsymbol{<\,1\%}$ \\
\midrule
One C384 forecast job memory footprint       & hundreds of GB
                                                across many compute nodes \\
One PR sweep's compute-node-hours            & $10^{3}\text{--}10^{4}$
                                                core-hours \\
Master-tick CPU work saved                   & \textbf{0}
                                                (same total, just serialised) \\
\bottomrule
\end{tabular}
\end{table}

So as a fraction of either (a) the scron service node or (b) the
resources the PR sweep itself consumes downstream on compute nodes,
the change sits in the noise floor --- well under~1\%.

\subsection{Where the Bound Is Not Negligible}

The one regime where the change is meaningful is when the
\code{cgroup} applied to the scron job is intentionally small:

\begin{itemize}[leftmargin=*]
  \item If the cron partition allocates each scron job e.g.\
        \textbf{1~core / 2~GB} (a common defensible default), then
        $M_\text{peak} = 0.5\text{--}3$\,GB \emph{straddles or exceeds the
        limit}, causing OOM-kill of the tick. This is almost
        certainly the failure mode the Gaea CM observed; the PR fixes
        it cleanly.
  \item If the cgroup is generous ($\geq 8$\,GB, $\geq 4$\,cores),
        the change is essentially cosmetic.
\end{itemize}

\subsection{Aggregate Site-Wide Impact}

If every active GW developer simultaneously ran a 10-case PR sweep
(say $\sim 5$ concurrent users), the aggregate difference in Slurm
controller load is on the order of:

\begin{equation}
\Delta \;\approx\;
  \underbrace{5}_{\text{users}} \cdot
  \underbrace{10}_{\text{exp/user}} \cdot
  \frac{1\,\text{burst}}{300\,\text{s}}
  \;\approx\; \mathbf{0.17\,\text{rpc/s}}
\end{equation}

against a controller that comfortably handles
$10^{3}\text{--}10^{4}$\,rpc/s. Invisible.

\section{Throughput Trade-off}
\label{sec:throughput}

The actual model jobs' wallclock is unchanged --- Rocoto still submits
them to Slurm exactly as before. What changes is the cadence at which
each experiment's task graph is re-evaluated:

\begin{itemize}[leftmargin=*]
  \item \textbf{Before:} every experiment's \code{rocotorun} fires
        every 5\,min in parallel; per-experiment polling cadence is
        constant in $N$.
  \item \textbf{After:} the master walks experiments serially within a
        single 5-min tick. If the cumulative pass time
        ($N\cdot t$) stays well below $T$, cadence is preserved. If
        it does not, \code{singleton} blocks the next tick and
        per-experiment cadence stretches roughly to
        $N\cdot t / T \cdot 5$\,min.
\end{itemize}

This is bounded but real: it scales linearly in $N$ and in
\code{rocotorun}'s pass time. For current PR loads the effect is
sub-second to sub-minute and immaterial; at higher case counts or
with large XMLs it would degrade end-to-end PR-sweep wallclock by
adding idle time between a task completing and the next dependent
task being submitted. The current implementation provides no
mitigation (e.g.\ capping the master at~$<5$\,min, or sharding into
2--3 master jobs).

\section{Recommendations}
\label{sec:recommendations}

\begin{enumerate}[leftmargin=*]
  \item \textbf{Reframe the rationale.} The accurate one-line
        description is: \emph{``prevent OOM-kill of the
        scron-managed Rocoto tick on Gaea by serialising per-experiment
        passes,''} not an aggregate efficiency improvement.
  \item \textbf{Validate the \code{cgroup} assumption} on Gaea (and any
        other scron-using site) so the trade-off is documented and
        site-portable.
  \item \textbf{Per-site partition/account.} The master block reuses
        the partition/account of the \emph{first} experiment's
        crontab; this is fine when all PR cases target the same
        service partition (typical) but would silently misroute the
        driver tick if a future case ever specified a different one.
        Consider asserting equality across collected
        \code{*.crontab} files before emitting the master block.
  \item \textbf{Bound the master pass.} Add a guard inside
        \code{rocoto\_master\_run.sh} so that if the cumulative
        wallclock approaches~$T$, remaining experiments are deferred
        to the next tick rather than relying solely on
        \code{--dependency=singleton} to absorb the overrun.
\end{enumerate}

\section{Conclusion}
\label{sec:conclusion}

The branch correctly addresses a real, site-specific operational
failure on Gaea --- OOM-kill of the cron-driven Rocoto tick under a
constrained \code{cgroup} --- using a clean, contained change to a
single CM utility. Its framing as a memory or efficiency improvement
overstates the system-level impact: as a fraction of either the scron
node or the PR sweep's compute footprint, the saving is below~1\%,
and it introduces a bounded, $O(N)$ throughput cost in worst-case
polling cadence. The trade is favourable for the present PR sweep
size; it would warrant revisiting if PR case counts or
\code{rocotorun} pass time grow materially.

\end{document}
